\documentclass[12pt]{article}
\usepackage[spanish]{babel}
\usepackage[utf8]{inputenc}
\usepackage{geometry}
\geometry{margin=2.5cm}

\title{Arquitecturas Distribuidas en la Industria: Caso Netflix}
\author{Jeremy Alexis Alvarez Parraga}
\date{}

\begin{document}

\maketitle

\section{Caso de Estudio}

Netflix es una de las empresas más conocidas por la implementación de una arquitectura basada en microservicios. Inicialmente utilizaba una arquitectura monolítica, pero debido al crecimiento acelerado de usuarios y servicios, migró a una arquitectura distribuida compuesta por cientos de microservicios independientes.

\section{Razones para Elegir Microservicios}

La empresa adoptó este estilo arquitectónico para mejorar la escalabilidad, disponibilidad y mantenimiento de la plataforma. Cada servicio puede desarrollarse, desplegarse y actualizarse de forma independiente, permitiendo responder rápidamente a las necesidades del negocio.

\section{Beneficios Obtenidos}

Entre los principales beneficios destacan:

\begin{itemize}
\item Escalabilidad independiente de cada servicio.
\item Mayor tolerancia a fallos.
\item Despliegues continuos sin afectar toda la plataforma.
\item Equipos de desarrollo más autónomos.
\item Mejor experiencia para millones de usuarios simultáneos.
\end{itemize}

\section{Desafíos Encontrados}

La adopción de microservicios también presentó desafíos importantes:

\begin{itemize}
\item Mayor complejidad en la comunicación entre servicios.
\item Necesidad de herramientas de monitoreo y observabilidad.
\item Gestión de la consistencia de datos distribuidos.
\item Incremento en los costos de infraestructura.
\end{itemize}

\section{Aplicación al Proyecto Propio}

En el proyecto desarrollado durante la asignatura, la arquitectura de microservicios podría aplicarse separando funcionalidades como autenticación, gestión de usuarios, procesamiento de datos y generación de reportes. Esto permitiría que cada módulo evolucione de manera independiente, facilitando el mantenimiento y la escalabilidad futura del sistema.

\section{Conclusión}

La experiencia de Netflix demuestra que los microservicios constituyen una alternativa efectiva para sistemas que requieren alta disponibilidad y crecimiento continuo. Aunque su implementación implica desafíos técnicos, los beneficios obtenidos en términos de flexibilidad y escalabilidad justifican su adopción en aplicaciones distribuidas modernas.

\begin{thebibliography}{9}

\bibitem{hasselbring}
Hasselbring, W., \& Steinacker, G. (2017).
Microservice Architectures for Scalability, Agility and Reliability in E-Commerce.
IEEE International Conference on Software Architecture Workshops (ICSAW).
DOI: 10.1109/ICSAW.2017.11

\bibitem{pahl}
Pahl, C., \& Jamshidi, P. (2016).
Microservices: A Systematic Mapping Study.
Proceedings of the 6th International Conference on Cloud Computing and Services Science.
DOI: 10.5220/0005785501370146

\end{thebibliography}
\end{document}
